% ============================================================================
% ЗАКЛЮЧЕНИЕ, без номер
% ============================================================================
\section*{ЗАКЛЮЧЕНИЕ}
\addcontentsline{toc}{section}{ЗАКЛЮЧЕНИЕ}

В курсовата работа е разгледано възстановяването на изтрити и повредени видеофайлове
като подготвителен етап на видеоанализа в криминалистиката. Изводите са следните.

Първо, възможностите за възстановяване се определят от структурата на формата. Във
файл MP4 видеоданните в \code{mdat} не носят сами информация за границите на кадрите. Тази
информация, както и параметрите на кодека, са в кутията \code{moov}. Затова загубата на
\code{moov} прави файла нечетим, дори когато всички кадри са налице.

Второ, всяка причина за загуба оставя различни следи на носителя. Изтриването запазва
данните, форматирането губи връзката им с файловете, прекъснатият запис губи
метаданните, фрагментирането разпокъсва файла, а презаписването унищожава част от
данните безвъзвратно.

Трето, на тези случаи съответстват различни методи: възстановяване по файловата
система, изрязване по сигнатури, сглобяване на фрагменти с валидиране, възстановяване
на ниво кадри и ремонт с еталонен файл. Колкото по-малко е оцеляло от структурата на
файла, толкова повече методът разчита на знание за формата и толкова повече от
резултата е изчислено, а не прочетено от носителя.

Четвърто, доказателствената стойност на резултата зависи от процедурата. Работата върху
копие, контролните суми, веригата на съхранение и ясното разграничаване между
възстановено и реконструирано съдържание са това, което позволява възстановеният запис
да бъде използван в разследването.

Работата е ограничена до формата MP4 и кодека H.264. Естествено продължение е да се
разгледат други контейнери и кодеци, както и проверката за автентичност на
възстановения запис.
